% =============================================================
% Chapitre 3 - Analyse des besoins et modélisation des menaces (~9 pages)
% Source : ../../CH3_Analyse_des_besoins_et_modele_de_menaces.md
% =============================================================
% -- Resserrement typographique local au chapitre 3 (valeurs restaurées en fin de fichier).
\setlength{\textfloatsep}{10pt plus 2pt minus 3pt}
\setlength{\floatsep}{8pt plus 2pt minus 2pt}
\setlength{\intextsep}{8pt plus 2pt minus 2pt}
\setlength{\abovecaptionskip}{4pt}
\setlength{\belowcaptionskip}{2pt}
\setlength{\LTpre}{6pt}\setlength{\LTpost}{6pt}

\chapter{Analyse des besoins et modélisation des menaces}
\label{ch:besoins-menaces}
\soustitrechapitre{De la valeur métier à l'exigence de sécurité vérifiable}

Une architecture de sécurité ne se conçoit pas à partir d'une liste de bonnes pratiques, mais à
partir d'une réponse à trois questions : qu'est-ce qui a de la valeur, qui voudrait s'y attaquer,
et par quel chemin ? Ce chapitre y répond, puis en dérive vingt exigences de sécurité et la matrice
qui relie le risque analysé ici, l'architecture du \cref{ch:conception}, la réalisation du \cref{ch:realisation} et les
mesures du \cref{ch:validation}.

\textbf{Pourquoi deux méthodes plutôt qu'une.} Aucune des deux méthodes retenues au \cref{ch:etat-art} ne
suffit seule : EBIOS Risk Manager hiérarchise et fixe la gravité, mais laisserait des scénarios
trop abstraits pour en dériver un test ; STRIDE énumère systématiquement, mais produirait un
inventaire sans hiérarchie de valeur, où le balayage automatisé pèserait autant que la perte du
corpus. Leur combinaison sépare le niveau métier, qui décide de ce qui compte, du niveau technique,
qui décide d'où cela peut être atteint.

\section{Acteurs et cas d'utilisation}

\subsection{Acteurs légitimes et acteurs hostiles}

L'analyse distingue les acteurs \textbf{légitimes}, dont les besoins définissent les
fonctionnalités, des acteurs \textbf{hostiles}, dont les objectifs définissent les menaces. La
distinction n'est pas étanche --- un compte légitime détourné devient un acteur hostile disposant
de droits légitimes --- d'où une série unique A1 à A11 dont une colonne indique la nature
(\cref{tab:acteurs}), plutôt qu'une numérotation séparée qui suggérerait deux populations
disjointes.

{\footnotesize
\setlength{\tabcolsep}{3pt}
\begin{longtable}{|P{2.5cm}|P{1.3cm}|P{5.4cm}|P{5.85cm}|}
\caption[Acteurs du modèle, A1 à A11]{Acteurs du modèle : série unique, nature indiquée en colonne}
\label{tab:acteurs} \\
\hline
\textbf{Acteur} & \textbf{Nature} & \textbf{Rôle ou motivation} & \textbf{Interaction ou moyens} \\
\hline
\endfirsthead
\caption[]{Acteurs du modèle (suite)} \\
\hline
\textbf{Acteur} & \textbf{Nature} & \textbf{Rôle ou motivation} & \textbf{Interaction ou moyens} \\
\hline
\endhead
A1 — Utilisateur & Légitime & Contribue à l'application hébergée & Point d'entrée public (F1) \\ \hline
A2 — Analyste & Légitime & Consulte, qualifie et documente les incidents & Tableau de bord, lecture et verdict \\ \hline
A3 — Administrateur & Légitime & Administre la plateforme & Consoles d'administration et dépôts \\ \hline
A4 — Développeur & Légitime & Modifie le code ou l'infrastructure & Chaînes F3 et F7 \\ \hline
A5 — Identité de déploiement & Légitime & Identité automatisée à privilèges & Fédération sans secret permanent, F3 et F7 \\ \hline
A6 — Tâches internes & Légitime & Détectent et enrichissent & Requêtes et tâches planifiées, F5 \\ \hline
A7 — Opportuniste & Hostile & Gain rapide, cible non choisie & Outils automatisés, vulnérabilités connues \\ \hline
A8 — Acteur intéressé par le corpus & Hostile & Valeur du corpus linguistique constitué & Ciblage délibéré, patience \\ \hline
A9 — Participant fraudeur & Hostile & Gain financier lié aux prix de la compétition & Comptes multiples, manipulation applicative \\ \hline
A10 — Chaîne d'approvisionnement & Hostile & Exécution de code chez la victime & Dépendance ou image de base compromise \\ \hline
A11 — Interne ou prestataire & Hostile & Négligence, curiosité ou malveillance & Accès légitime initial, connaissance du système \\
\hline
\end{longtable}
}

Deux acteurs appellent une vigilance particulière. L'identité de déploiement (A5) possède des
droits de modification sur l'ensemble du socle et doit être traitée comme une identité à
privilèges, non comme un outil. Le participant fraudeur (A9) dispose d'un accès parfaitement
légitime et d'une incitation financière : aucun contrôle de périmètre ne le distingue d'un
contributeur ordinaire, et c'est le comportement applicatif, non le réseau, qui le révèle.

\subsection{Cas d'utilisation du socle}

Huit cas d'utilisation structurent le socle, formulés du point de vue de l'acteur et non du
composant technique (\cref{fig:cas-utilisation}), chacun rattaché à un flux : F1 pour UC1 à
UC3, F3 pour UC4, F7 pour UC5, F4 et F5 pour UC6, F5 pour UC7, F6 pour UC8. Deux commandent la
conception : \textbf{UC3} impose que le verdict de l'analyste soit enregistré sans pouvoir altérer
les preuves sur lesquelles il porte, d'où la séparation, au \cref{ch:conception}, entre tables de preuve en
écriture interdite et table de verdict humain ; \textbf{UC8}, qui correspond à la question Q6 et au
verrou V4, constitue l'apport propre du mémoire.

\begin{figure}[htbp]
\centering
\begin{tikzpicture}[scale=0.92, every node/.append style={transform shape},
  uc/.style={draw=black, thick, ellipse, inner sep=1.5pt, text width=2.65cm,
    minimum width=3.9cm, minimum height=1.05cm, align=center, font=\scriptsize,
    fill=black!4},
  acteur/.style={font=\scriptsize, align=center, text width=1.7cm},
  acteurSysteme/.style={draw=black, thick, rounded corners=2pt, fill=black!6,
    font=\scriptsize, align=center, text width=1.65cm, minimum height=0.8cm},
  humain/.style={-, thick, draw=black!75, rounded corners=2pt},
  machine/.style={-, thick, dashed, draw=black!55, rounded corners=2pt}
]
  % Frontière du système. Les gouttières intérieures (haut, bas, gauche, droite)
  % et le couloir central entre les deux colonnes restent libres : aucune
  % association ne longe le cadre ni n'en sort pour y rentrer.
  \draw[thick, rounded corners=5pt, fill=black!1] (1.55,0.20) rectangle (12.00,8.85);
  \node[font=\small\bfseries, fill=white, inner sep=2pt] at (3.40,8.85)
    {Socle Zero Trust};

  \node[uc] (uc1) at (4.35,7.15) {UC1\\Accéder à une application};
  \node[uc] (uc2) at (4.35,5.45) {UC2\\Consulter les alertes};
  \node[uc] (uc3) at (4.35,3.75) {UC3\\Qualifier une alerte\\et enregistrer le verdict};
  \node[uc] (uc5) at (4.35,1.95) {UC5\\Modifier l'infrastructure};
  \node[uc] (uc4) at (9.05,7.15) {UC4\\Livrer une version};
  \node[uc] (uc6) at (9.05,5.45) {UC6\\Détecter un comportement\\suspect};
  \node[uc] (uc7) at (9.05,3.75) {UC7\\Enrichir une alerte};
  \node[uc] (uc8) at (9.05,1.95) {UC8\\Prioriser les correctifs};

  % Acteurs humains à gauche.
  \foreach \y/\id/\nom in {7.15/A1/Utilisateur final,4.60/A2/Analyste sécurité,1.95/A3/Administrateur}{
    \draw[thick] (0.55,\y+0.36) circle (0.14) (0.55,\y+0.22) -- (0.55,\y-0.24)
      (0.32,\y) -- (0.78,\y) (0.55,\y-0.24) -- (0.33,\y-0.60)
      (0.55,\y-0.24) -- (0.77,\y-0.60);
    \node[acteur, anchor=north] at (0.55,\y-0.68) {\textbf{\id}\\\nom};
  }
  % Le développeur est humain ; l'identité de déploiement et les tâches
  % planifiées sont des systèmes.
  \draw[thick] (13.55,7.51) circle (0.14) (13.55,7.37) -- (13.55,6.91)
    (13.32,7.15) -- (13.78,7.15) (13.55,6.91) -- (13.33,6.55)
    (13.55,6.91) -- (13.77,6.55);
  \node[acteur, anchor=north] at (13.55,6.47) {\textbf{A4}\\Développeur};
  % A5 descendu de 5,05 a 4,70 : a 5,05 sa bordure haute recouvrait le libelle
  % « Developpeur » de A4, ancre au nord en 6,47 et haut de deux lignes.
  \node[acteurSysteme] (a5) at (13.55,4.70) {\textbf{A5}\\\guillemotleft système\guillemotright\\Identité de déploiement};
  \node[acteurSysteme] (a6) at (13.55,2.60) {\textbf{A6}\\\guillemotleft système\guillemotright\\Tâches planifiées};

  % Associations humaines à gauche : ramification courte dans la gouttière.
  \draw[humain] (0.78,7.15) -- (uc1.west);
  \draw[humain] (0.78,4.60) -- (2.05,4.60);
  \draw[humain] (2.05,4.60) |- (uc2.west);
  \draw[humain] (2.05,4.60) |- (uc3.west);
  \draw[humain] (0.78,1.95) -- (uc5.west);

  % A4 : UC4 en vis-à-vis ; UC5 et UC8 par un tronc unique qui franchit la
  % gouttière haute puis descend dans le couloir libre entre les deux colonnes.
  \draw[humain] (13.32,7.15) -- (uc4.east);
  \draw[humain] (13.55,7.65) -- (13.55,8.15) -- (6.70,8.15) -- (6.70,1.95) -- (uc5.east);
  \draw[humain] (6.70,2.60) -- (uc8.140);

  % A5 et A6 : acteurs système ancrés face au bord droit du cadre, segment court
  % vers l'ellipse la plus proche. Le seul trajet long contourne la colonne
  % d'acteurs par la droite et rejoint UC5 par la gouttière basse.
  \draw[machine] (a5.north west) -- (uc4.338);
  \draw[machine] (a5.east) -- (14.80,5.05) -- (14.80,0.72) -- (4.35,0.72) -- (uc5.south);
  \draw[machine] (a6.north west) -- (uc6.332);
  \draw[machine] (a6.north west) -- (uc7.342);
  \draw[machine] (a6.west) -- (uc8.10);

  \draw[humain] (3.00,-0.35) -- +(0.7,0) node[right, font=\footnotesize] {association humaine};
  \draw[machine] (7.80,-0.35) -- +(0.7,0) node[right, font=\footnotesize] {association automatisée};
\end{tikzpicture}
\caption{Diagramme de cas d'utilisation du socle}
\label{fig:cas-utilisation}
\end{figure}

\section{Spécification des besoins}

\subsection{Besoins fonctionnels}

Les besoins fonctionnels traduisent les carences C1 à C4 identifiées par l'audit du \cref{ch:cadre-existant}
(\cref{tab:besoins-fonctionnels}).

{\footnotesize
\setlength{\tabcolsep}{3pt}
\begin{longtable}{|P{1cm}|P{9.85cm}|P{2.1cm}|P{2.1cm}|}
\caption{Besoins fonctionnels\label{tab:besoins-fonctionnels}} \\
\hline
\textbf{\#} & \textbf{Besoin} & \textbf{Segment} & \textbf{Origine} \\
\hline
\endfirsthead
\caption[]{Besoins fonctionnels (suite)} \\
\hline
\textbf{\#} & \textbf{Besoin} & \textbf{Segment} & \textbf{Origine} \\
\hline
\endhead
BF1 & Exposer les applications derrière un point d'entrée unique, filtré et journalisé & Sécurité & C1 \\ \hline
BF2 & Attribuer une identité distincte à chaque charge de travail, droits minimaux & Sécurité & C1 \\ \hline
BF3 & Rendre les données inaccessibles depuis Internet, accès privé uniquement & Sécurité & C1 \\ \hline
BF4 & Décrire l'intégralité du socle sous forme de code versionné & Ops & C2 \\ \hline
BF5 & Contrôler automatiquement chaque livraison avant mise en production & Dev & C3 \\ \hline
BF6 & Authentifier l'identité de déploiement sans secret permanent & Dev & C3 \\ \hline
BF7 & Centraliser les journaux de toutes les couches dans un entrepôt unique & Détection & C4 \\ \hline
BF8 & Exécuter des règles de détection versionnées sur ces journaux & Détection & C4 \\ \hline
BF9 & Garantir qu'aucun composant d'analyse ne modifie les preuves qu'il analyse & Sécurité & (propre) \\ \hline
BF10 & Associer automatiquement chaque alerte à une technique d'attaque connue & Données & V3 \\ \hline
BF11 & Présenter à l'analyste incidents, contexte et conduite à tenir & Métier & C4 \\ \hline
BF12 & Réordonner les vulnérabilités détectées à la livraison selon les menaces observées & Données & V4 \\ \hline
BF13 & Permettre l'accueil d'une nouvelle application sans reconception du socle & Métier & (entreprise) \\
\hline
\end{longtable}
}

\textbf{BF9} est volontairement absolu : \emph{un moteur de détection ne doit jamais modifier les
preuves qu'il analyse}. \textbf{BF13} traduit l'objectif d'entreprise du \cref{ch:cadre-existant} : aucun
composant du socle ne doit contenir de logique spécifique à une application.

\subsection{Besoins non fonctionnels}

Le \cref{tab:besoins-non-fonctionnels} recense les sept besoins non fonctionnels, avec
pour chacun la cible visée et le moyen de vérification retenu.

{\footnotesize
\setlength{\tabcolsep}{3pt}
\begin{longtable}{|P{1.2cm}|P{3.4cm}|P{7.65cm}|P{2.8cm}|}
\caption{Besoins non fonctionnels}
\label{tab:besoins-non-fonctionnels} \\
\hline
\textbf{\#} & \textbf{Besoin} & \textbf{Cible} & \textbf{Vérification} \\
\hline
\endfirsthead
\caption[]{Besoins non fonctionnels (suite)} \\
\hline
\textbf{\#} & \textbf{Besoin} & \textbf{Cible} & \textbf{Vérification} \\
\hline
\endhead
BNF1 & Disponibilité & Objectif de niveau de service instrumenté, mesuré en continu & Supervision, ch.6 \\ \hline
BNF2 & Latence & Objectif défini et mesuré ; tout dépassement documenté & Supervision, ch.6 \\ \hline
BNF3 & Coût mensuel & Quelques dizaines d'euros ; poste dominant identifié & Facturation, ch.6 \\ \hline
BNF4 & Exploitabilité par une seule personne & Procédures écrites et exécutées au moins une fois & Exécution réelle, ch.6 \\ \hline
BNF5 & Reproductibilité & Reconstruction complète d'un environnement, chronométrée & Reconstruction complète \\ \hline
BNF6 & Auditabilité & Chaque décision automatisée journalisée, entrées et version & Inspection des journaux \\ \hline
BNF7 & Scalabilité & Volontairement limitée ; limites mesurées et documentées & Déclaration explicite, ch.6 \\
\hline
\end{longtable}
}

\textbf{BNF3} : le poste de coût dominant d'une supervision de sécurité n'est pas le calcul mais
l'\textbf{ingestion des journaux}, d'où le choix de filtrer à la source (\cref{ch:conception}). \textbf{BNF7}
est une limite assumée : déclarée sans mesure, elle resterait une affirmation, d'où son
inscription au programme de validation du \cref{ch:validation}.

\section{Flux de données et frontières de confiance}
\label{sec:dfd-frontieres}

STRIDE ne s'applique pas à une architecture en général mais à un objet précis : un
\textbf{diagramme de flux de données} portant ses \textbf{frontières de confiance} --- les lieux
où une donnée passe sous le contrôle d'une entité de confiance différente. C'est là, et seulement
là, que la méthode énumère ses six catégories : poser ce diagramme est la condition de validité des
scénarios opérationnels qui en sont dérivés.

\subsection{Le diagramme de flux de données du socle}

La \cref{fig:dfd-frontieres} représente le socle dans la notation usuelle de la modélisation de
menaces. Les composants y portent exactement les noms de la vue de déploiement du \cref{ch:conception}
(\cref{fig:deploiement}), afin que les deux figures se lisent l'une sur l'autre ; les arcs portent
les identifiants des sept flux F1 à F7 (\cref{tab:flux}).

\begin{figure}[htbp]
\centering
\begin{tikzpicture}[
  x=0.925cm, y=0.90cm, font=\scriptsize,
  proc/.style={draw=black!70, rounded corners=3pt, fill=black!4,
               align=center, inner sep=2.5pt, minimum height=0.62cm},
  mag/.style={draw=black!70, double, double distance=0.7pt, fill=white,
              align=center, inner sep=2.5pt, minimum height=0.58cm},
  ext/.style={draw=black!70, fill=black!10, align=center,
              inner sep=2.5pt, minimum height=0.58cm},
  cadre/.style={draw=black!35, densely dotted, rounded corners=4pt},
  tb/.style={draw=black!80, dashed, line width=0.9pt},
  tblab/.style={fill=white, inner sep=1.5pt, font=\scriptsize\bfseries},
  a/.style={-{Latex[length=1.5mm]}, black!65},
  ab/.style={{Latex[length=1.5mm]}-{Latex[length=1.5mm]}, black!65},
  ai/.style={-{Latex[length=1.5mm]}, black!55, densely dotted},
  et/.style={font=\scriptsize, fill=white, inner sep=1pt}]

% ============ frontieres empilees TB1..TB4, franchies par le reseau ============
% Libelles volontairement reduits au code : les cinq frontieres sont enoncees
% en toutes lettres dans le tableau qui suit immediatement la figure.
\draw[tb] (0,10.95) -- (9.55,10.95);   \node[tblab] at (0.32,11.14) {TB1};
\draw[tb] (0,9.15)  -- (9.55,9.15);    \node[tblab] at (0.32,9.34)  {TB2};
\draw[tb] (0,7.25)  -- (9.55,7.25);    \node[tblab] at (0.32,7.44)  {TB3};
\draw[tb] (0,4.75)  -- (9.55,4.75);    \node[tblab] at (0.32,4.94)  {TB4};

% ============ frontiere orthogonale TB5, franchie par l'identite seule ========
\draw[tb] (9.95,2.55) -- (9.95,12.95);
\node[tblab] at (9.95,13.16) {TB5};

% ============ entites externes, au-dessus de TB1 ==============================
\node[ext, text width=2.0cm]  (uti)   at (1.15,12.15) {Utilisateurs de\\l'application hébergée};
\node[ext, text width=1.8cm]  (ana)   at (4.10,12.15) {Analyste et\\administrateur};
\node[ext, text width=2.2cm]  (forge) at (7.55,12.15) {Dépôt de code et\\chaîne de livraison};
\node[ext, text width=3.5cm]  (ref)   at (12.95,12.15)
      {Référentiel de techniques d'attaque,\\catalogues de vulnérabilités};

% ============ bande 1 : perimetre =============================================
\node[proc, text width=8.2cm] (per) at (4.75,10.15)
      {\textbf{Périmètre (L1)} \; --- \; répartiteur HTTPS, filtrage applicatif, limitation de débit};

% ============ bande 2 : charges de travail ====================================
\node[proc, text width=2.15cm] (app) at (1.45,8.20) {Application\\hébergée};
\node[proc, text width=2.15cm] (tdb) at (4.75,8.20) {Tableau de\\bord};
\node[proc, text width=2.15cm] (api) at (8.05,8.20) {API de\\supervision};

% ============ bande 3 : reseau prive ==========================================
\node[cadre, minimum width=8.9cm, minimum height=1.30cm] (rp) at (4.75,6.15) {};
\node[anchor=north west, font=\scriptsize\itshape, black!55] at (0.45,6.78) {Réseau privé (L4)};
\node[proc, text width=2.8cm] (enc) at (2.35,5.98) {Service d'encodage\\(entrée interne seule)};
\node[proc, text width=2.8cm] (enr) at (7.15,5.98) {Tâche planifiée\\d'enrichissement};

% ============ bande 4 : donnees ===============================================
\node[mag, text width=3.2cm] (bdd) at (3.10,3.85) {Base de données};

% ============ plan manage du fournisseur, a droite de TB5 =====================
\node[mag,  text width=3.4cm] (reg) at (12.95,10.60) {Registre d'images\\et coffre de secrets};
\node[proc, text width=3.4cm]  (col)  at (12.95,8.90)  {Collecte journalisée (L7)};
\node[mag,  text width=3.4cm]  (siem) at (12.95,7.25)  {Entrepôt de supervision};
\node[proc, text width=3.4cm]  (req)  at (12.95,5.65)  {Requêtes planifiées\\de détection};
\node[mag,  text width=3.4cm]  (obj)  at (12.95,4.05)  {Stockage objet};

% ============ arcs ============================================================
% -- TB1 : entree publique
\draw[a] (uti) -- node[et, pos=0.30] {F1} (uti|-per.north);
\draw[a] (ana) -- (ana|-per.north);
% -- TB2 : perimetre vers charges de travail
\draw[a] (per.south-|app) -- node[et, pos=0.45] {F1} (app.north);
\draw[a] (per.south-|tdb) -- (tdb.north);
\draw[a] (per.south-|api) -- (api.north);
% -- le tableau de bord n'atteint les donnees que par l'API
\draw[a] (tdb.east) -- node[et, pos=0.5] {jeton} (api.west);
% -- TB3 puis TB4 : acces aux donnees par le reseau prive
\draw[ab] (app.south) -- node[et, pos=0.62, xshift=-3.2mm] {F2} (app|-bdd.north);
\draw[ab] (api.south) .. controls +(0,-1.1) and +(2.4,0.2) .. (bdd.east);
% -- TB5 : l'API atteint l'entrepot par l'identite seule, hors reseau prive
\draw[a] (api.east) -- node[et, pos=0.52] {\textbf{F2b} --- identité seule (É11)} (siem.west);
% -- collecte des journaux
\draw[a] (per.east) -- node[et, pos=0.40] {F4} (col.west);
\draw[a] (col.south) -- (siem.north);
% -- enrichissement semantique
\draw[a] (enr.west) -- node[et, pos=0.5] {F5a} (enc.east);
\draw[ab] (enr.east) -- node[et, pos=0.55] {F5b} (siem.south west);
% -- detection : entierement dans le plan manage
\draw[ab] (req.north) -- (siem.south);
% -- chaine de livraison : identite federee, sans cle
\draw[a] (forge.east) -- node[et, pos=0.45] {F3} (reg.north west);
% F7 descend en ligne droite a x=10,35 : il doit etre trace AVANT l'etiquette
% « etiquette SHA », dont le fond blanc le masque a l'endroit ou ils se croisent.
% Trace apres, il barrait le S et le donnait a lire comme un signe dollar.
  \draw[a] (forge.south east) -- (10.35,11.50) -- (10.35,4.05) -- node[et, pos=0.93] {F7} (obj.west);
\draw[ai] (reg.south) -- ++(0,-0.42) -- node[et, pos=0.78] {étiquette SHA} (api.north east);
% -- imports differes de referentiels
\draw[ai] (ref.south) .. controls +(1.9,-1.1) and +(1.9,1.3) .. node[et, pos=0.55] {F6} (siem.east);

% ============ legende et cartouche ============================================
\node[draw=black!45, rounded corners=2pt, fill=black!2, align=left,
      inner sep=4pt, text width=6.0cm, font=\scriptsize] at (3.35,1.50)
  {\textbf{Notation} \\[1pt]
   \tikz[baseline=-0.5ex]{\node[proc, minimum width=0.5cm, minimum height=0.22cm, inner sep=1pt]{};}\;processus \quad
   \tikz[baseline=-0.5ex]{\node[mag, minimum width=0.5cm, minimum height=0.22cm, inner sep=1pt]{};}\;magasin de données \\
   \tikz[baseline=-0.5ex]{\node[ext, minimum width=0.5cm, minimum height=0.22cm, inner sep=1pt]{};}\;entité externe \quad
   \tikz[baseline=-0.5ex]{\draw[tb] (0,0) -- (0.65,0);}\;frontière de confiance \\
   \tikz[baseline=-0.5ex]{\draw[ai] (0,0) -- (0.65,0);}\;flux différé ou hors requête \\[1pt]
   \emph{TB1 à TB5 sont énoncées en toutes lettres au tableau suivant.}};


\end{tikzpicture}
\caption[Diagramme de flux de données du socle et ses cinq frontières de confiance]{Diagramme
de flux de données du socle et ses cinq frontières de confiance. Les frontières TB1 à TB4 sont
empilées et franchies par le réseau ; TB5 leur est orthogonale et n'est franchie que par
l'identité. C'est cette frontière, sur laquelle aucun dispositif réseau ne peut s'interposer,
que traversent douze des vingt-quatre scénarios opérationnels, dont six des huit
scénarios de gravité critique.}
\label{fig:dfd-frontieres}
\end{figure}

{\footnotesize
\setlength{\tabcolsep}{2pt}
\begin{longtable}{|P{0.75cm}|P{2.45cm}|P{2.2cm}|P{2.35cm}|P{3.6cm}|P{3.5cm}|}
\caption[Les cinq frontières de confiance et les scénarios qui les franchissent]{Les cinq
frontières de confiance : ce que chacune porte, et les scénarios opérationnels qui la franchissent}
\label{tab:frontieres-scenarios} \\
\hline
\textbf{\#} & \textbf{Frontière} & \textbf{Franchie par} & \textbf{Menaces dominantes} & \textbf{Nature du contrôle qui s'y interpose} & \textbf{Scénarios rattachés (nombre, dont critiques)} \\
\hline
\endfirsthead
\caption[]{Les cinq frontières de confiance (suite)} \\
\hline
\textbf{\#} & \textbf{Frontière} & \textbf{Franchie par} & \textbf{Menaces dominantes} & \textbf{Nature du contrôle qui s'y interpose} & \textbf{Scénarios rattachés (nombre, dont critiques)} \\
\hline
\endhead
TB1 & Internet $\rightarrow$ périmètre & Requête HTTPS (F1) & Déni de service, divulgation par reconnaissance, usurpation & Réseau et applicatif : filtrage applicatif, limitation de débit, terminaison TLS & SO2 à SO5 --- \textbf{4}, aucun critique \\ \hline
TB2 & Périmètre $\rightarrow$ charges de travail & Trafic déjà filtré (F1) & Altération, élévation de privilège & Réseau : entrée des services restreinte au répartiteur ; jeton d'identité vérifié & SO1 --- \textbf{1}, aucun critique \\ \hline
TB3 & Charges de travail $\rightarrow$ réseau privé & Interface de sortie directe de la charge (F2, F5) & Divulgation par canal sortant, élévation par rebond & Réseau : refus par défaut en sortie, autorisations nommées par identité et par port, refus journalisés & SO17, SO18, SO24 --- \textbf{3}, aucun critique \\ \hline
TB4 & Réseau privé $\rightarrow$ données & Adresse privée (F2) & Divulgation d'information & Réseau : absence d'adresse publique, chiffrement obligatoire du transport & SO6, SO8 --- \textbf{2}, les deux critiques \\ \hline
TB5 & Charges de travail $\rightarrow$ \textbf{plan managé du fournisseur} & \textbf{Identité seule}, hors du réseau privé (F3, F4, F5b, F7) & Usurpation, altération des preuves, élévation de privilège & \textbf{Identité exclusivement} : autorisations, jetons de courte durée, fédération & SO7, SO9 à SO14, SO16, SO20 à SO23 --- \textbf{12}, dont \textbf{6 critiques} : SO9, SO10, SO12, SO20, SO21, SO22 \\
\hline
\end{longtable}
}

\subsection{TB5, frontière sans contrôle réseau}

Les quatre premières frontières se comportent comme celles d'une architecture classique : chacune
sépare deux zones que le trafic relie par un chemin réseau, sur lequel un dispositif peut être
posé — filtre, refus par défaut, absence d'adresse publique. Un lecteur habitué au modèle
périmétrique les reconnaît, et c'est ce qui rend TB5 difficile à voir.

\textbf{TB5 n'est pas un chemin réseau.} Lorsqu'une charge de travail écrit dans l'entrepôt de
supervision, lit un secret dans le coffre, tire une image du registre ou modifie l'état de
l'infrastructure, elle s'adresse aux interfaces du fournisseur : aucun dispositif réseau
n'arbitre cet appel. Depuis la refonte du plan réseau (\cref{subsec:reseau-cible}), il emprunte
bien l'interface de sortie de la charge et se heurte au refus par défaut ; mais la règle qui le
laisse passer ne désigne qu'une plage de destination du fournisseur, la même pour l'entrepôt du
socle et pour tout autre entrepôt du même fournisseur. Il est autorisé, ou refusé, par la seule
vérification de l'identité qui le présente. Le \cref{ch:conception} relève ce point comme l'écart É11 : la
vue de déploiement, qui aligne les composants sur un plan réseau, laisse croire l'inverse.

Deux conséquences en découlent. Aucun contrôle réseau ne réduit un risque porté par TB5 : aucun ne
sépare l'interface du fournisseur que le socle emploie de celle qu'un attaquant contrôlerait chez
le même fournisseur, et les distinguer supposerait un périmètre de service, donc une organisation
dont le projet ne dispose pas. Et TB5 est la frontière que franchissent les chemins
d'administration et de livraison --- précisément ceux que le modèle périmétrique du \cref{ch:cadre-existant}
laissait sans contrôle parce qu'il ne les voyait pas comme des chemins.

\begin{keybox}
Sur TB5, l'identité n'est pas un contrôle parmi d'autres : elle est le seul. C'est la traduction
opérationnelle du principe d'identité (PR2) et la raison pour laquelle le \cref{ch:conception} traite le
plan d'identité comme la segmentation réelle de l'architecture.
\end{keybox}

\section{Analyse de risque}

\subsection{Valeurs métier, sources de risque et scénarios stratégiques}

Une \textbf{valeur métier} est ce dont la perte, l'altération ou la divulgation nuit à
l'organisation ; un \textbf{bien support} est un élément technique dont une valeur métier dépend.
Six valeurs métier ont été retenues : le corpus linguistique constitué, premier actif de
l'entreprise ; les identités et les voix des contributeurs ; l'intégrité de la compétition, dont
dépend la crédibilité du dispositif d'incitation ; la disponibilité du service ; l'intégrité des
preuves de supervision ; la crédibilité technique de MENAL, atteinte transverse des cinq
précédentes. La cinquième est propre au socle : protéger les données sans protéger les traces de
l'attaque rend la supervision contestable, et c'est elle qui justifie l'interdiction d'écriture
faite aux composants d'analyse (BF9).

Neuf biens supports les portent, de la base applicative au point d'entrée public, et les cinq
sources de risque sont les acteurs hostiles A7 à A11 du \cref{tab:acteurs}. Une sixième,
l'acteur étatique organisé, est \textbf{écartée par décision explicite} : la traiter supposerait
des moyens de détection et de réponse sans commune mesure avec le périmètre du projet, et l'inclure
sans ces moyens produirait une analyse dont aucune exigence ne découlerait ; l'exclusion est datée
et révisable si le corpus change de statut. La codification des ateliers figure en
annexe~\ref{ann:modele-menaces}.

\textbf{Sept scénarios stratégiques} relient chacun une source de risque à une valeur métier
atteinte, par un chemin de haut niveau ; l'annexe~\ref{ann:modele-menaces} les énonce avec leur
codification et leur gravité. \textbf{Quatre sont de gravité critique} : accès direct à la base de
données puis copie du corpus et des identités, obtention du jeton d'export du corpus, détournement
de l'identité de déploiement suivi d'un déploiement non revu, divulgation d'un secret permettant de
reconstituer des accès légitimes. Les trois autres, graves ou significatifs, portent sur
l'altération des journaux et des détections, sur le classement de la compétition et sur la
disponibilité d'un service exposé. Trois des quatre critiques --- accès direct à la base, jeton
d'export, divulgation d'un secret --- correspondent à des points bloquants de l'audit du
\cref{ch:cadre-existant} : l'audit a fourni la matière factuelle de cette analyse, non une liste théorique.

\subsection{Scénarios opérationnels}

Les scénarios stratégiques sont traduits en chemins techniques par application systématique des
six catégories STRIDE aux sept flux F1 à F7, en chacun des points où ils franchissent une des cinq
frontières de la \cref{sec:dfd-frontieres}. L'opération produit \textbf{vingt-quatre
scénarios opérationnels SO1 à SO24}, dont huit au niveau de risque le plus élevé ; le tableau
détaillé de l'annexe~\ref{ann:modele-menaces} rattache chacun à sa frontière et à sa catégorie
STRIDE, et la dernière colonne du \cref{tab:frontieres-scenarios} en donne la répartition.

\textbf{Lecture.} Douze scénarios sur vingt-quatre, et \textbf{six des huit de gravité critique},
franchissent TB5 : ce résultat n'était pas un objectif de la modélisation, il en sort. Les deux
scénarios de TB4 sont eux aussi de gravité critique, mais un contrôle réseau simple et vérifiable
les adresse --- l'absence d'adresse publique.

Trois rattachements méritent d'être justifiés, aucun n'ayant été forcé. \textbf{SO15}, l'absence
de trace permettant d'attribuer une action, et \textbf{SO19}, la falsification du rattachement
d'une vulnérabilité pour la déclasser, ne franchissent aucune frontière et ne figurent donc dans
aucune ligne du tableau ; ce sont aussi les seuls auxquels la matrice des techniques d'attaque ne
fait correspondre aucune technique, pour la même raison --- un défaut de traçabilité et un abus de
droit interne ne sont pas des chemins d'attaque, ce sont les conditions qui les rendent
indétectables. \textbf{SO1} est rattaché à TB2, la frontière qu'une charge utile franchit
effectivement, alors que le contrôle qui le traite est posé une frontière en amont, sur TB1 : seul
endroit du socle où la défense est placée avant la frontière menacée plutôt que sur elle.

\subsection{Appréciation et traitement du risque}

Le croisement gravité $\times$ vraisemblance, détaillé en annexe~\ref{ann:modele-menaces}, place
la zone de risque la plus élevée sur les chaînes de livraison applicative et de provisionnement
(F3, F7) et sur l'accès direct aux données (F2) --- la répartition qui précède en donne
l'explication structurelle. Le scénario le plus critique de vraisemblance forte correspond
exactement au point bloquant B2 de l'audit.

Les décisions de traitement sont, pour l'essentiel, la \textbf{réduction} par des contrôles
d'architecture, complétée par deux \textbf{acceptations explicites}, datées et révisables ; aucun
transfert n'est appliqué, faute d'assurance ou de sous-traitance de sécurité dans le périmètre.

\section{Exigences de sécurité et traçabilité}

\subsection{Les vingt exigences de sécurité}

L'analyse de risque est traduite en vingt exigences de sécurité EX1 à EX20
(\cref{tab:exigences}). Chacune est formulée comme une \textbf{propriété vérifiable du
système} et non comme une intention, ce qui rend possible l'écriture d'un test au \cref{ch:validation} :
« les accès doivent être maîtrisés » ne se teste pas ; « une connexion directe depuis Internet
est refusée » se teste par une connexion directe depuis Internet.

{\footnotesize
\setlength{\tabcolsep}{3pt}
\begin{longtable}{|P{1cm}|P{14.5cm}|}
\caption{Formulation complète des vingt exigences de sécurité}
\label{tab:exigences} \\
\hline
\textbf{\#} & \textbf{Exigence de sécurité} \\
\hline
\endfirsthead
\caption[]{Formulation complète des vingt exigences de sécurité (suite)} \\
\hline
\textbf{\#} & \textbf{Exigence de sécurité} \\
\hline
\endhead
EX1 & Toute requête empruntant le point d'entrée public est soumise aux règles de filtrage configurées \\ \hline
EX2 & Les tentatives d'authentification anormales sont limitées et détectées \\ \hline
EX3 & Le service reste disponible sous charge anormale et le balayage est tracé \\ \hline
EX4 & La base applicative n'expose pas d'adresse publique et une connexion directe depuis Internet est refusée \\ \hline
EX5 & Une identité applicative ne peut accéder qu'aux données de son périmètre \\ \hline
EX6 & Tout accès aux données est lié à une identité vérifiée, sans jeton transmissible \\ \hline
EX7 & Les comptes de service du périmètre n'utilisent pas de clé statique inventoriée ; toute exception est tracée \\ \hline
EX8 & Un artefact présentant une vulnérabilité critique corrigeable détectée par la porte configurée est bloqué \\ \hline
EX9 & Le déploiement ne peut tirer que du registre de confiance unique \\ \hline
EX10 & Un secret reconnu par les règles du détecteur configuré bloque la chaîne avant livraison \\ \hline
EX11 & Les identités d'exécution des composants d'analyse n'ont aucun droit d'écriture sur les tables de preuve \\ \hline
EX12 & Toute interruption de la collecte de journaux est détectable \\ \hline
EX13 & Toute action privilégiée est attribuable à un acteur identifié \\ \hline
EX14 & Le service d'encodage n'accepte que des appels internes, et sa sortie réseau est restreinte aux destinations nécessaires à son fonctionnement \\ \hline
EX15 & Une charge anormale d'alertes ne compromet pas la chaîne d'enrichissement \\ \hline
EX16 & Une modification hors code d'une ressource déjà gérée est détectée par rafraîchissement ; les ressources non gérées sont contrôlées par inventaire \\ \hline
EX17 & L'état de l'infrastructure est privé, versionné et chiffré \\ \hline
EX18 & Toute modification de droits est revue avant application \\ \hline
EX19 & Une consommation de ressources anormale est détectée \\ \hline
EX20 & Les sorties réseau sont limitées à des destinations connues, et les refus du pare-feu sont journalisés \\
\hline
\end{longtable}
}

Deux exigences appellent une précision d'état à la date de \dateref{}, postérieure à la refonte du
plan réseau (décision D15, \cref{subsec:reseau-cible}, qui en établit le détail et les dates).
\textbf{EX20} est portée par un contrôle et non plus par l'autorisation implicite du fournisseur :
la totalité du trafic sortant des charges porteuses de données emprunte le réseau virtuel, cinq
autorisations nommées par identité et par port y sont posées, et le refus par défaut en sortie
rejette tout le reste ; les destinations autorisées sont celles inventoriées pendant la fenêtre du
15 au 22/08/2026, non celles supposées, et les refus sont journalisés depuis le 19/08/2026 (écart
É1, clos). Le critère se vérifie en deux points : aucune règle d'autorisation ne comporte de plage
source, et une destination hors inventaire est refusée puis retrouvée dans l'entrepôt --- objet du
protocole T20.

\textbf{EX14} est portée en entrée par la restriction du service d'encodage aux appels internes,
et contre le rebond par son retrait du réseau privé : le composant ne porte plus d'interface dans
un sous-réseau, et son identité dédiée --- écart É3, corrigé le 19/08/2026 --- ne détient que le
droit d'écrire des journaux. Sa sortie emprunte en revanche le chemin managé de la plateforme
d'exécution, qu'aucune règle de pare-feu ne gouverne ; le relevé du 25/08/2026 l'établit, et
l'exigence reste ouverte à ce titre (protocole T14, non conforme).

\textbf{Deux réserves accompagnent ces états}, toutes deux établies en
\cref{subsec:reseau-cible} : l'exfiltration par les interfaces du fournisseur n'est pas fermée, et
le sélecteur des règles est l'étiquette réseau et non le compte de service --- la frontière est
donc opposable à une charge de travail compromise, non à un opérateur capable de déployer. La
distinction entre la propriété exigée et le degré de garantie atteint est portée par le
\cref{ch:validation}, jamais dissimulée dans la formulation de l'exigence.

\subsection{Matrice de traçabilité}

Cette matrice (\cref{tab:matrice-tracabilite}) relie les quatre niveaux d'analyse : pour
chaque scénario opérationnel, elle établit la chaîne complète du risque analysé jusqu'à sa
démonstration. Chaque ligne se lit ainsi : \emph{ce scénario impose cette exigence, réalisée par
ce contrôle, dont l'efficacité est démontrée par ce test.} L'exigence y est appelée par son
identifiant ; sa formulation complète est celle du \cref{tab:exigences} ci-dessus.

{\footnotesize
\setlength{\tabcolsep}{3pt}
\begin{longtable}{|P{1.5cm}|P{1.0cm}|P{6.3cm}|P{6.2cm}|}
\caption[Matrice de traçabilité menace--exigence--contrôle--test]{Matrice de traçabilité
menace $\rightarrow$ exigence $\rightarrow$ contrôle $\rightarrow$ test}
\label{tab:matrice-tracabilite} \\
\hline
\textbf{Menace} & \textbf{Exig.} & \textbf{Contrôle (couche)} & \textbf{Test de validation} \\
\hline
\endfirsthead
\hline
\textbf{Menace} & \textbf{Exig.} & \textbf{Contrôle (couche)} & \textbf{Test de validation} \\
\hline
\endhead
SO1 & EX1 & Filtrage applicatif au point d'entrée unique (L1) & T1 — Charges d'attaque connues ; vérification du blocage \\ \hline
SO2, SO3 & EX2 & Limitation de débit (L1) ; règle de détection sur échecs répétés (L5) & T2 — Échecs répétés ; blocage et alerte \\ \hline
SO4, SO5 & EX3 & Limitation de débit, restriction géographique (L1) ; journalisation (L7) & T3 — Test de charge ; limitation et trace \\ \hline
SO6 & EX4 & Base sans adresse publique, accès privé uniquement (L4, L5) & T4 — Connexion directe externe ; échec attendu \\ \hline
SO7 & EX5 & Une identité de service par charge de travail (L2) & T5 — Lecture de l'entrepôt de supervision ; refus attendu \\ \hline
SO8 & EX6 & Authentification requise entre services (L2, L3) & T6 — Appel sans jeton ; refus attendu \\ \hline
SO9 & EX7 & Identité fédérée, jetons courts, création de clés interdite (L2) & T7 — Inventaire des clés ; tentative de création ; refus attendu \\ \hline
SO10 & EX8 & Contrôle de vulnérabilité bloquant (F3) & T8 — Artefact vulnérable ; blocage attendu \\ \hline
SO11 & EX9 & Registre unique autorisé, droits limités (L2, F3) & T9 — Déploiement depuis source externe ; refus attendu \\ \hline
SO12 & EX10 & Détection de secrets bloquante (F3) ; gestionnaire dédié (L2) & T10 — Secret factice ; blocage attendu \\ \hline
SO13, SO16 & EX11 & Droits d'écriture réservés à la collecte (L2, L5) & T11 — Écriture avec l'identité d'enrichissement ; refus attendu \\ \hline
SO14 & EX12 & Supervision de la continuité du flux (L7) & T12 — Interruption volontaire ; alerte attendue \\ \hline
SO15 & EX13 & Journaux d'audit du fournisseur activés (L7) & T13 — Action privilégiée ; attribution vérifiée \\ \hline
SO17 & EX14 & Entrée interne (L6) ; composant retiré du réseau privé, sans interface ni chemin vers une ressource privée ; identité dédiée limitée à l'écriture de journaux (L2, L4) & T14 — Connexion sortante émise depuis le composant ; refus et trace attendus \\ \hline
SO18 & EX15 & Traitement par lots à cadence fixe (L6) & T15 — Volume anormal ; comportement et coût observés \\ \hline
SO20 & EX16 & Infrastructure en code exclusive, détection de dérive (F7) & T16 — Ressource manuelle ; détection comme dérive \\ \hline
SO21 & EX17 & Stockage d'état distant, accès restreint (F7) & T17 — Accès non autorisé ; refus attendu \\ \hline
SO22 & EX18 & Approbation manuelle obligatoire (F7) & T18 — Changement élargissant des droits ; blocage en attente \\ \hline
SO23 & EX19 & Supervision de consommation, alerte de budget (L7) & T19 — Consommation simulée ; alerte attendue \\ \hline
SO24 & EX20 & Refus par défaut en sortie ; cinq autorisations nommées par identité et par port ; refus journalisés et versés à l'entrepôt (L4, É1) & T20 — Destination hors inventaire ; refus attendu et ligne de refus retrouvée dans l'entrepôt \\
\hline
\end{longtable}
}

\textbf{Vingt exigences, vingt tests.} Ces tests constituent le programme de validation du
\cref{ch:validation}, qu'ils fixent \emph{avant} que celui-ci ne soit écrit ; aucun n'est optionnel, une
exigence ne devenant une propriété établie qu'au passage de son test. La matrice se lit dans les
deux sens : de gauche à droite elle justifie chaque exigence par le scénario qui la motive, de
droite à gauche elle vérifie que chacune porte son test. Elle compte une ligne par exigence
et nomme donc vingt des vingt-quatre scénarios ; le traitement des quatre autres --- SO3,
SO5, SO16 et SO19 --- est arrêté par la table de décision de risque de l'\cref{ann:modele-menaces},
qui les couvre tous. Elle établit enfin que les douze scénarios
de TB5 sont adressés par des contrôles d'identité (EX5 à EX11, EX16 à EX19) et non par des
contrôles réseau qui ne s'y appliqueraient pas.

\section{Rattachement aux référentiels}

Les vingt exigences sont rattachées aux référentiels du \cref{ch:etat-art}
(\cref{tab:rattachement-referentiels}), pour exprimer la conformité en termes reconnus et
identifier les écarts assumés ; le \cref{ch:validation} en donne le détail, lacunes comprises.

{\footnotesize
\setlength{\tabcolsep}{3pt}
\begin{longtable}{|P{3.4cm}|P{4.4cm}|P{7.5cm}|}
\caption{Rattachement des exigences aux référentiels}
\label{tab:rattachement-referentiels} \\
\hline
\textbf{Référentiel} & \textbf{Exigences concernées} & \textbf{Écarts assumés} \\
\hline
\endfirsthead
\caption[]{Rattachement des exigences aux référentiels (suite)} \\
\hline
\textbf{Référentiel} & \textbf{Exigences concernées} & \textbf{Écarts assumés} \\
\hline
\endhead
Zero Trust (\cite{nist800207}) & EX4-EX7, EX11, EX13, EX14, EX20 & Aucun principe non traité ; degré mesuré au ch.6 \\ \hline
Applicatif (\cite{owasp2025}) & EX1, EX8-EX10 & Code hébergé hors périmètre \\ \hline
Infrastructure (CIS) & EX4, EX5, EX7, EX13, EX16, EX17, EX20 & Journalisation des flux échantillonnée en régime permanent, justifiée par le coût \\ \hline
Techniques d'attaque (\cite{mitre_attack}) & Couverture des scénarios opérationnels & Partielle par construction ; plafond et lacunes publiés au ch.6 \\
\hline
\end{longtable}
}

\section*{Conclusion du chapitre}
\addcontentsline{toc}{section}{Conclusion du chapitre}

Ce chapitre a transformé les carences du \cref{ch:cadre-existant} et l'état de l'art du \cref{ch:etat-art} en exigences
vérifiables : six valeurs métier, dont l'intégrité de la chaîne de preuve, propre au socle ; cinq
sources de risque, la sixième écartée par décision datée ; sept scénarios stratégiques, que
l'application systématique de STRIDE au diagramme de flux et à ses cinq frontières a traduits en
vingt-quatre scénarios opérationnels, puis en vingt exigences associées chacune à un contrôle et à
un test. Le principal enseignement est double : les scénarios les plus graves ne concernent pas
l'application exposée sur Internet mais les chemins d'administration, de livraison et d'accès aux
données, et ces chemins se distinguent par le fait qu'ils franchissent TB5, la seule frontière du
socle qu'aucun contrôle réseau ne peut garder. Le chapitre suivant présente l'architecture qui
réalise ces vingt contrôles, et son plan d'identité --- réponse directe à TB5.

% -- Restauration des valeurs par défaut d'espacement des flottants.
\setlength{\textfloatsep}{20pt plus 2pt minus 4pt}
\setlength{\floatsep}{12pt plus 2pt minus 2pt}
\setlength{\intextsep}{12pt plus 2pt minus 2pt}
\setlength{\abovecaptionskip}{10pt}
\setlength{\belowcaptionskip}{0pt}
\setlength{\LTpre}{\bigskipamount}\setlength{\LTpost}{\bigskipamount}
